\documentclass[pdflatex, sn-nature]{sn-jnl}

%%%% Standard Packages
%%<additional latex packages if required can be included here>

\usepackage{graphicx}%
\usepackage{multirow}%
\usepackage{amsmath, amssymb, amsfonts, cancel}%
\usepackage{amsthm}%
\usepackage[title]{appendix}%
\usepackage{xcolor}%
\usepackage{textcomp}%
\usepackage{manyfoot}%
\usepackage{booktabs}%
\usepackage{algorithm}%
\usepackage{algorithmicx}%
\usepackage{algpseudocode}%
\usepackage{listings}%
\usepackage{enumitem}%
\usepackage{tikz}%
\usepackage[colorinlistoftodos]{todonotes}
\usetikzlibrary{bayesnet,positioning,fit,calc}%

\theoremstyle{thmstyleone}%
\newtheorem{theorem}{Theorem}%
\newtheorem{proposition}[theorem]{Proposition}%

\theoremstyle{thmstyletwo}%
\newtheorem{example}{Example}%
\newtheorem{remark}{Remark}%

\theoremstyle{thmstylethree}%
\newtheorem{definition}{Definition}%

% insert comments
\setlength{\marginparwidth }{2cm}
\usepackage{todonotes}

% multicol environements
\usepackage{multicol}

%%%% maths operators
% Sets
\newcommand \NN {\mathbb{N}}
\newcommand \ZZ {\mathbb{Z}}
\newcommand \QQ {\mathbb{Q}}
\newcommand \RR {\mathbb{R}}

% Indicator
\newcommand \indic {\mathbbm{1}}

% Probability Operators
\DeclareMathOperator*{\prob}{\mathbb{P}}
\DeclareMathOperator*{\esp}{\mathbb{E}}
\DeclareMathOperator*{\var}{\mathbb{V}\text{ar}}
\DeclareMathOperator*{\cov}{\mathbb{C}\text{ov}}
\newcommand \PP [1]{\prob\left({#1}\right)}
\newcommand \EE [1]{\esp\left[{#1}\right]}
\newcommand \VV [1]{\var\left[{#1}\right]}
\newcommand \CC [1]{\cov\left[{#1}\right]}
% Law
\newcommand \unif [1] {\mathcal{U}\left({#1}\right)}
\newcommand \normal [2] {\mathcal{N}\left({#1},{#2}\right)}

% Operators
\DeclareMathOperator*{\argmax}{arg\,max}
\DeclareMathOperator*{\argmin}{arg\,min}
% Matrix
\DeclareMathOperator*{\diag}{Diag}
\DeclareMathOperator*{\rank}{Rank}
\DeclareMathOperator*{\DET}{Det}
\DeclareMathOperator*{\Tr}{Tr}
\DeclareMathOperator*{\ES}{ES} 
\DeclareMathOperator*{\COV}{Cov}
\DeclareMathOperator*{\Spectrum}{Spectrum}

% To name R packages
\newcommand{\CRANpkg}[1]{\href{https://CRAN.R-project.org/package=#1}{\textbf{#1}}}
\newcommand{\strong}[1]{\texorpdfstring%
{{\normalfont\fontseries{b}\selectfont #1}}%
{#1}}
\newcommand{\code}[1]{\verb|#1|}

% Colours of the single-cell derivation (Appendix): technical scales and composition.
\definecolor{scalecol}{HTML}{C45C26}
\definecolor{compcol}{HTML}{1D4E89}


\raggedbottom

\begin{document}
    \title[DeCovarT: A Probabilistic and Multidimensional Framework for Cellular
    Deconvolution in Heterogeneous Biological Samples]{DeCovarT: A Probabilistic
    and Multidimensional Framework for Cellular Deconvolution in Heterogeneous Biological
    Samples}

    \author*[1,2]{\fnm{Bastien} \sur{Chassagnol}}
    \email{bastien\_chassagnol@laposte.net}
    \equalcont{B.C.\ and M.P.\ contributed equally to this work.}

    \author[3]{\fnm{Marielle} \sur{P\'{e}r\'{e}}}
    \email{marielle.pere@inria.fr}
    \equalcont{B.C.\ and M.P.\ contributed equally to this work.}

    \author[1]{\fnm{Mohamad} \sur{Al Charif}}

    \author[2]{\fnm{Gr\'{e}gory} \sur{Nuel}}
    \email{nuel@math.cnrs.fr}

    \author[4]{\fnm{Etienne} \sur{Becht}}
    \email{etienne.becht@inserm.fr}

    \author[1]{\fnm{Ana\"{i}s} \sur{Baudot}}
    \email{anais.baudot@univ-amu.fr}

    \affil[1]{\orgdiv{Marseille Medical Genetics (MMG), Systems Biomedicine team},
      \orgname{Aix-Marseille University, Inserm U1251},
      \orgaddress{\street{Facult\'{e} de M\'{e}decine, 27 Bd Jean Moulin},
        \city{Marseille}, \postcode{13385}, \country{France}}}

    \affil[2]{\orgdiv{Laboratoire de Probabilit\'{e}s, Statistique et
      Mod\'{e}lisation (LPSM), UMR 8001},
      \orgname{Sorbonne Universit\'{e} / CNRS},
      \orgaddress{\city{Paris}, \country{France}}}

    \affil[3]{\orgdiv{MacBes team},
      \orgname{Inria Centre at Universit\'{e} C\^{o}te d'Azur;
        Institut de Pharmacologie Mol\'{e}culaire et Cellulaire (IPMC)},
      \orgaddress{\city{Sophia-Antipolis}, \country{France}}}

    \affil[4]{\orgdiv{Centre de Recherche sur l'Inflammation (CRI),
      UMR 1149},
      \orgname{Inserm / Universit\'{e} Paris Cit\'{e}},
      \orgaddress{\street{Facult\'{e} de m\'{e}decine -- site Bichat,
        16 rue Henri Huchard},
        \city{Paris}, \postcode{75018}, \country{France}}}


        \abstract{Although bulk transcriptomic analyses have greatly contributed to
        a better understanding of complex diseases, their sensibility is hampered by
        the highly heterogeneous cellular compositions of biological samples. To
        address this limitation, computational deconvolution methods have been designed
        to automatically estimate the frequencies of the cellular components that make
        up tissues, typically using reference samples of physically purified
        populations. However, they perform badly at differentiating closely related
        cell populations.
    
        We hypothesised that the integration of the covariance matrices of the
        reference samples could improve the performance of deconvolution algorithms.
        We therefore developed a new tool, DeCovarT, that integrates the structure of
        individual cellular transcriptomic network to reconstruct the bulk profile.
        Specifically, we inferred the ratios of the mixture components by a standard
        maximum likelihood estimation (MLE) method, using the Levenberg-Marquardt
        algorithm to recover the maximum from the parametric convolutional
        distribution of our model. We then consider a reparametrisation of the log-likelihood
        to explicitly incorporate the simplex constraint on the ratios. Preliminary
        numerical simulations suggest that this new algorithm outperforms previously
        published methods, particularly when individual cellular transcriptomic profiles
        strongly overlap.
        }

        \keywords{deconvolution, bulk RNA-seq, weighted multivariate Gaussian distribution, covariance-driven model}

    \maketitle


    \section{Introduction}

   \subsection{context: Bulk versus single cell analysis, and application to organoid develoment and embryogenes}
    \label{sec1} Bulk versus single cell analysis
    The analysis of the bulk transcriptome provided new insights on
    the mechanisms underlying disease development. However, such methods ignore the
    intrinsic cellular heterogeneity of complex biological samples, by averaging
    measurements over several distinct cell populations. Failure to account for changes
    of the cell composition is likely to result in a loss of \emph{specificity}
    (genes mistakenly identified as differentially expressed, while they only
    reflect an increase in the cell population naturally producing them) and \emph{sensibility}
    (genes expressed by minor cell populations are amenable being masked by highly
    variable expression from major cell populations).
    
   - cell-type composition: a criticla cue in early devleopment stages to predict the moprhotypic fate of gastruloids

   - On the othar hand, single cell analysis provides a more detailed picture of the cellular heterogeneity,
   enabling to idenity with an uncuperendted elvel of resolution cell ypes and cell states 

   - However, their cost and the elated computaitonal burden prevent shteir use for large-scale analysis.


    Accordingly, a range of computational methods have been developed to
    estimate cellular fractions, but they perform poorly in discriminating cell types
    displaying high phenotypic proximity. 
    
    \subsection{Research gap: evlauate the untertainty associated with the estimated cell-type proportions and
     discrmiant closely related cell types with simialr averaged expresison profiles}
    
    Indeed, most of them assume that purified
    cell expression profiles are fixed observations, omitting the variability
    and intrinsically interconnected structure of the transcriptome. For instance,
    the gold-standard deconvolution algorithm \emph{CIBERSORT} \cite{newmanRobustEnumerationCell2015}
    applies nu-support vector regression ($\nu$-SVR) to recover the minimal subset
    of the most informative genes in the purified signature matrix. However, this
    machine learning approach assumes that the transcriptomic expressions are independent. 

    - relying on regression-based frameworks or deep elanring approaches prevent from 
    deriving confidnce intervals on the estimated cell-type proportions. Accoridngly, msot approaches only provide point estimates.

    

       \todo{Furether detail related approaches to set apart closley related cell types, such as RNA-Sieve which incrproates cell type specific variability and
     imapct of the scaling factor induced by different sequencing tehcnoloiges but assuming no interactions between genes. MEAD uses a error-in-model based apporach
     , but cna be rather consider as a theorerical paper trying to relax a smuch as posisble assumptions to derive  asymoptotic confidence intervals when the statistical distirbution mapping biolgoicla 
     and tehcnical sources of variability with total bulk expression is unknown. Finally, Decals mplements a cosntrained (to enforce the unti-simpelx constraint) generalised elast squares ap^poroach
     by iteratively inferring the cellualr ratios and the cell type speciifc covariances. Hence, in cosntrast wiht our aporiahc, it jointly infers these two critica aprameters, 
     which we claim is liley to lead in a pooly idneitable problem. }


     \subsection{DeCovarT: a covariance-driven framework for bulk deconvolution}

     In contrast to these approaches, we hypothesised that derivign a covariance-driven 
     estimation of the cell-type proportions by integrating cell-type specific (CTS) covariance profiles
      could enhance the performance of transcriptomic deconvolution methods.
      \emph{DeCovarT} (Covariance-AwareDeconvolution of Bulk transcritpomics), is, the first generative probabilistic model that releis on a parametrised
      \emph{multivariate} Gaussian convolution for bulk deconvolution. 
    
    - Ohter innovations: an extented synthetic generative model describing a variety of scenariii and versatile enought to arametrise
    differencet covariance structures. Levle of shannon entropy (to describe the cellar heterogeinity), targeted cosine simialrityof the eman expression profile
    using a Gram-based approach, and most improtantly an overlooked metric, the overlap, that has a nice itnerpretion as a n estimate of the probability to missassign
    a cell type base don its osberved expression profile, even when the true estimates are available -> it hus provides alower bound on the untertainty of the estimated cell-type proportions.

    \todo{Idea of the overlap is to demosntrate that evne in a simpel clasisfication case, we faila t a certain rpobability levle, there's no way of
    recovering the composition of that clel type in the arder deconvolution setting framework}

    
    demosntrating the role on ideal cases of accoutnign for the covariance, not only to derive proper cofnidence itnervals 
    with a higher coverage  but returning least bieased estimates
    - 
    \subsection{Layout and directio of the paper: the workflow for building (and potentiallye xtending) DeCovarT}

    - in ..., we describe our approach for incorproating synthetic netowkr
    - in ..., our framewokr for building relasitic and versatile scenariir testing different structures of correlation between genes
    - in ..., we describe one approach to tailor decvart's,dieal model to the specificty of single-cell, toa cocutn for disparties betwene cross-paltofrm
    (and suggest in ... other extended approaches)

    \section{Results}
    \subsection{Overview of the method}
    - an overiew of th eenrative model and the pipleine is porpsoed in Figure 1
    - comapred to oter deconvolution apporoaches, we i,corproate the cell-type speicfic variability and demosnrate that conseirng
    the vaeraged expression alone is not enougfh to both select markers of interest, and 
    tends to systeetically bias the estiamtion of closely related cell types. The storngest imapt is not
    on the point wise estimation( as quantifid by the bias), but rather  on the coverage of confidenc eitnervals
    - we also demosntrate that our approach enables to spearate closley relatd clel types, notbly
    if they share similr averaged expression profiles but are are charzcthered by strong entokr discrpenacies. We hwoever
    demosntrate theoretically using the widfth of the confdinence itnervals a s a prox that it's usualy ahrder to infer rare clel types, 
    implying to resort to lternative apporaches or colelcting them as a residual mixtur enosie term; 
   
    \subsection{Synthetic scenarios}

    \subsubsection{The toy example with two genes and two cell types}

    \subsubsection{A more realist and higher dimensional framework}

    \subsection{Application on real biological use cases: organoid development and embryogenesis}
    
     - gastrulod sue case: major limitation preventing the dmeocratistion of organoids is the high variability
     of viability (and phjortype profiles) across batches, even when devleopped n simialr condtions
     - besides, contrary to more mature cell tyes, sucha s the differntiated immune cellscirculatin in the blood, 
     develoomental cels usually lie on a continuum of differentiation with gradient potency, rather than than fixed discrete states, rather than being sharply defined.
     - comparison of decovart here with other deocnvolution apporoaches, and the limitations of the current state of the art
    \section{Methods}

    \subsection{Notations}

    \subsection{Model assumptiosn to derive the enrative model}

    \todo{Add an illustration of the model: 
      top panel: the generative model, 
      bottom panel: the pipelien to apply DeCovarT on real world example, from the idtification of informative genes to delenate closely related cell types to the derivatio of confidence intervals
      including comptation of the scaling ratio and derivation of a sparse structure of the precision matri}

    \subsection{Derivation of the covariance-driven deconvolution ideal model}

    - derive an ideal model without accoutnng foo across-platform diffrentiated trancriptomic capture efficienty

    \subsection{Inference strategy to recover the amximum likelihood estimates of DeCovarT's model}

    \subsection{How to generat interesting synthetic scnearios with an epxoential comiantion of hyper-parameters to be tested}

    - desription of statstical metrics to quantify the cmomplexity of the reocnsutrctuin of the clualr composition
    - description fot he actual pipeline to generate synthetic datasets with targetted descriptive metrics

    \todo{Panel A: the different stages for generating synthetic datasets with targetted descriptive metrics,
    Panel B: illustration for interpeeting metrics}

    \subsection{Application on real biological use cases: organoid development and embryogenesis}

  \subsubsection{Brief description of the Suppinger dataset}
  
  \subsubsection{Omnibus, variance-aware approach for idneityfing genes able to discmrinate closely related cell types}
  - set apart two cases
  \begin{itemize}
    \item whe we have a signle biolgocla replciate
    \item when we have seeral biological replciaes
  \end{itemize}
    \subsubsection{How to tailor DeCovarT to the specificty of single-cell data}
    - Brief review of strategies to acoutn for disicnt sequenci dpeehts across sequencign tehcnologies and dpeending on the normaoisation protocl (by 
    rying to recover the absolute numer of cells per cell type, by averaging both the single-cell and bulk transcriptomic expessio pe gene, 
    or by considering normalised count per lbrary size,with stratgies such as TPM orRPKM)

    - briefly advocate against the absolute estimation oft he number of cells, and against the lirbayr normaisations rategy
    inducing to sovle a purely compositional framework. Major caveat: by dividng by the total number of observed trnscripts (the so-called lirbayr size),
    we induce a spurirous correlation between genesexpression due to the compsotional aporoach (inclduign for gesn athat were not rpeviously correlated)
    
    - Two caveats (not specific to DeCovarT):
      - we retrieve the relative transriptomic ontribtuion epr cell type, not the actual cell type compostion. 
      Two quanties are imiar if and only if both the transcriptomic extraction is simialr across cell types, 
      ad if each cell tpe has the same amount of absolute transcriptomic expresison -> both hyptheses lilkly not to hold

    
    
    \section{Perspectives}
     \subsection{Inference of interactions between genes is still an open-soruce and un-solved problem}

    A key practical challenge for \emph{DeCovarT} is that, unlike methods relying
    solely on mean expression profiles, it requires a covariance (or precision)
    matrix encoding gene--gene interactions.  Recovering this structure from
    observational expression data alone is inherently difficult: expression
    correlations are largely correlative, and inferring directed, mechanistically
    interpretable regulatory edges typically requires additional data modalities
    such as chromatin accessibility (ATAC-seq) or transcription-factor binding
    information \cite{badia-i-mompelGeneRegulatoryNetwork2023,maizelsGeneRegulatoryNetworks2026}.
    Systematic benchmarking further shows that even with richer multi-omic inputs,
    GRN inference can fail in common regimes \cite{nourisaGeneRNIBLivingBenchmark2025}.
    Whether recovering the true causal regulatory architecture is strictly
    necessary for accurate deconvolution, or whether statistically consistent
    covariance estimation is sufficient to improve compositional recovery, remains
    an open question.
